\begin{remark}
\label{remark-triangle}
We sketch an alternative, perhaps simpler, proof of the existence of
the fundamental triangle.
Let $A \to B \to C$ be ring maps and assume that $B \to C$ is injective.
Let $P_\bullet \to B$ be the standard resolution of $B$ over $A$ and
let $Q_\bullet \to C$ be the standard resolution of $C$ over $B$.
Picture
$$
\xymatrix{
P_\bullet : &
A[A[A[B]]] \ar[d]
\ar@<2ex>[r]
\ar@<0ex>[r]
\ar@<-2ex>[r]
&
A[A[B]] \ar[d]
\ar@<1ex>[r]
\ar@<-1ex>[r]
\ar@<1ex>[l]
\ar@<-1ex>[l]
&
A[B] \ar[d] \ar@<0ex>[l] \ar[r] &
B \\
Q_\bullet : &
A[A[A[C]]]
\ar@<2ex>[r]
\ar@<0ex>[r]
\ar@<-2ex>[r]
&
A[A[C]]
\ar@<1ex>[r]
\ar@<-1ex>[r]
\ar@<1ex>[l]
\ar@<-1ex>[l]
&
A[C] \ar@<0ex>[l] \ar[r] &
C
}
$$
Observe that since $B \to C$ is injective, the ring $Q_n$ is a
polynomial algebra over $P_n$ for all $n$. Hence we obtain a cosimplicial
object in $\mathcal{C}_{C/B/A}$ (beware reversal arrows).
Now set $\overline{Q}_\bullet = Q_\bullet \otimes_{P_\bullet} B$.
The key to the proof of Proposition \ref{proposition-triangle}
is to show that $\overline{Q}_\bullet$ is a resolution of $C$ over $B$.
This follows from Cohomology on Sites, Lemma
\ref{sites-cohomology-lemma-O-homology-qis}
applied to $\mathcal{C} = \Delta$, $\mathcal{O} = P_\bullet$,
$\mathcal{O}' = B$, and $\mathcal{F} = Q_\bullet$ (this uses that $Q_n$
is flat over $P_n$; see Cohomology on Sites, Remark
\ref{sites-cohomology-remark-simplicial-modules} to relate simplicial modules
to sheaves). The key fact implies that the distinguished triangle of
Proposition \ref{proposition-triangle}
is the distinguished triangle associated to the short exact sequence
of simplicial $C$-modules
$$
0 \to
\Omega_{P_\bullet/A} \otimes_{P_\bullet} C \to
\Omega_{Q_\bullet/A} \otimes_{Q_\bullet} C \to
\Omega_{\overline{Q}_\bullet/B} \otimes_{\overline{Q}_\bullet} C \to 0
$$
which is deduced from the short exact sequences
$0 \to \Omega_{P_n/A} \otimes_{P_n} Q_n \to \Omega_{Q_n/A} \to
\Omega_{Q_n/P_n} \to 0$ of
Algebra, Lemma \ref{algebra-lemma-ses-formally-smooth}.
Namely, by Remark \ref{remark-resolution} and the key fact the complex on the
right hand side represents $L_{C/B}$ in $D(C)$.

\medskip\noindent
If $B \to C$ is not injective, then we can use the above to get a
fundamental triangle for $A \to B \to B \times C$. Since
$L_{B \times C/B} \to L_{B/B} \oplus L_{C/B}$ and
$L_{B \times C/A} \to L_{B/A} \oplus L_{C/A}$
are quasi-isomorphism in $D(B \times C)$
(Lemma \ref{lemma-cotangent-complex-product})
this induces the desired distinguished triangle in $D(C)$
by tensoring with the flat ring map $B \times C \to C$.
\end{remark}